\section{Evaluation and reproducibility}
\label{sec:evaluation}

We evaluate the work at three levels: replay of the formal theorem, identity of
the deployed program, and one finalized Solana execution.  The associated
artifacts are published with the source repository
\cite{AspisRepository2026}.

\subsection{Formal replay}

The final accepted-execution theorem has a dependency closure of 331 tracked
Lean modules.  The release replay starts from a clean checkout, verifies that
every required source is tracked, applies the pinned Rust translator, rejects
unfinished proof markers and unchecked native evaluation, and builds the
entire closure with Lean~4.32.  It then prints the theorem's axioms.  The
recorded result contains Lean's standard logical foundations and no admitted
project theorem.

The replay is intentionally narrower than building every historical file in
the repository.  Its manifest is generated from the final theorem's actual
imports, so passing the replay establishes that the published dependency
closure is complete and self-contained.

\subsection{Program identity}

The deployment archive contains the clean source tree, lock file, build
command, Solana platform tools, Rust toolchain, program binary, proof body,
public statement, and signed transaction bytes.  A clean build with the
recorded tools reproduced the 1,258,496-byte program exactly.  SHA-256 digests
for every artifact and the complete source revision are listed in the
artifact guide.

This record gives a checkable chain
\[
  \text{archived source}
  \longrightarrow \text{program bytes}
  \longrightarrow \text{program identifier}
  \longrightarrow \text{transaction}.
\]
The first arrow is checked by a deterministic rebuild and byte comparison.
The second and third are checked against the captured loader and transaction
records.

\subsection{Mainnet execution}

The archived
\href{https://explorer.solana.com/tx/EJviPgF12i9iK2CveVaQSMeFQqDMFPQ1iPRUYEwNQE3zGquTUZNJXPZEENorcQtsnQj1orFmH1TPsgdbR3vJ2fE}
{mainnet transaction} finalized at slot 435,019,536.  It read the sealed proof
account, verified the 75,358-byte proof, and recorded the nullifier and new
pool state in 1,334,452 compute units.  The transaction limit was 1,356,912,
leaving 22,460 units of headroom.  Exact signed-wire simulation reported the
same cost.

The before-and-after account snapshots record the old root and sequence, the
new root and sequence, and creation of the nullifier marker.  A later cleanup
transaction closed the retained proof account and refunded 525,660,961
lamports while preserving the spend state.  The deployment was subsequently
closed, so the release archive includes the RPC responses and their digests
needed to inspect the historical accounts.

The demonstration pool contained one note.  This execution therefore measures
proof verification, compute cost, and the atomic state transition.  Privacy is
addressed by the distributional model of \cref{sec:privacy}, not by the size of
this demonstration's anonymity set.

\begin{table}[t]
\centering
\small
\caption{Recorded evaluation results.}
\label{tab:evaluation-results}
\begin{tabularx}{\textwidth}{@{}R{0.34\textwidth}Y@{}}
\toprule
Measurement & Result \\
\midrule
Formal dependency closure & 331 tracked Lean modules \\
Replay environment & Lean~4.32 with pinned Rust-to-Lean translator \\
Program binary & 1,258,496 bytes; clean build reproduced byte for byte \\
Proof body & 75,358 bytes \\
Public statement & 768 bytes \\
Finalized spend & slot 435,019,536 \\
Compute usage & 1,334,452 of 1,356,912 units \\
State effect & nullifier marker created; root and sequence updated atomically \\
Cleanup & 525,660,961 lamports refunded; spend state preserved \\
\bottomrule
\end{tabularx}
\end{table}

\subsection{How the evidence supports the claim}

Formal replay checks a universal theorem about the translated successful
execution.  Build reproduction identifies the program bytes to which the
source record refers.  The transaction supplies one concrete execution with a
measured cost and observed state change.  The conditional security theorem
then states the mathematical and cryptographic premises under which an
accepted false statement is bounded.  Each layer answers a different audit
question, and the artifact guide provides the exact command or digest for
checking it.
